\documentclass{article}
\usepackage{amsmath, amssymb}
\usepackage{pgfplots}

\begin{document}

\section{Introduction}

Today we will be covering basic definitions, equalities, and inequalities of information theory. Perhaps it is the most important lecture.

\section{Review}

$X$ is a discrete random variable (finite number of states) and entropy is defined as 
\[
H(X) = -\sum_a P(X=a)\log P(X=a)=-\mathbb{E}_X \log P(X)
\]
Shannon's entropy is the only function that satisfies the grouping axiom and 3 adition, veru reasonable, axioms.

If $X$ is a random variable over $n$ elements, then 
\[
0\leq H(X)\leq \log n 
\]

\subsection{Joint Entropy}

For random variables $X,Y$ then 
\[
H(X,Y) = -\mathbb{E}_{X,Y} \log P(X,Y)
\]
and it generalizes over any number of variables
\[
H(X_1,\ldots,X_n) = -\mathbb{E}_{X_1,\ldots,X_n} \log P(X_1,\ldots,X_n)
\]

\subsection{Conditional Entropy}

For random variables $X,Y$, then 
\[
H(X|Y) = \sum_a P(Y=b)\cdot H(X|Y=b)=- \mathbb{E}_{X,Y} \log P(X|Y)
\]
Intuitively, it is the uncertainty remains in $X$ that remains after seeing $Y$.

\section{Fundamental Equalities and Inequalities}

What can we say about the relations between:
\[
H(X),H(Y),H(X|Y)H(Y|X),H(X,Y)
\]
\subsection{Fundamental Inequality of Conditional Entropy}

For random variables $X,Y$
\[
0\leq H(X|Y)\leq H(X)
\]
The first inequality is trivial, and the second is proved by Jensen inequality.

Intuitively, the uncertainty remainded in $X$ after seeing $Y$ is smaller than the original uncertainty in $X$. The uncertainty can only decrease by revealing $Y$. Also the uncertainty is always larger than $0$.

$H(X|Y)=H(X)$ if and only if $X,Y$ are independent. It's not hard but we will not prove.

\subsubsection{Notable Example}

Note that for specific $b$, it could be the case the $H(X|Y=b)$ could be larger than $H(X)$. The following table (\ref{tab:joint_prob_dist}) is a good example.
\begin{table}[h]
\centering
\begin{tabular}{|c|c|c|}
\hline
& $X=0$ & $X=1$ \\ \hline
$Y=0 $ & $1/4$ & $1/4$  \\ \hline
$Y=1$ & $1/2$ & 0  \\ \hline
\end{tabular}
\caption{Joint Probability Distribution Table}
\label{tab:joint_prob_dist}
\end{table}

We have the following values
\begin{enumerate}
    \item $H(X|Y)=\frac{1}{2}$
    \item $H(X|Y=0)=1$
    \item $H(X)=H(\frac{3}{4},\frac{1}{4} )<1$
\end{enumerate}

\subsubsection{Proof of the Second Inequality}

See the photos.

\subsection{Chain Rule}

The most beautiful formula 
\[
H(X,Y) =H(Y) + H(X|Y)
\]
It's very intuitive! The uncertainty in $(X,Y)$ is the uncertainty in $Y$ plus the uncertainty remaining in $X$ after revealing $Y$.

It's a real treasure and is used many times in applications. 

We could have defined $H(X|Y)$ by this formula 
\[
H(X|Y) = H(X,Y) - H(Y)
\]
The chain rule is the most important property of the entropy function.

\subsubsection{Do we need to prove the chain rule?}

The chain rule is just the grouping axiom!

See the photo.

But we never proved the grouping axiom, so we do need to prove it.

\subsubsection{Proof of the Chain Rule}

\begin{align*}
    P(X,Y) &= P(Y)\cdot P(X|Y)\\
    \log P(X,Y) &= \log P(Y) + \log P(X|Y)\\
    \mathbb{E} \left[\log P(X,Y) \right]&= \mathbb{E}\left[\log P(Y)\right] + \mathbb{E}\left[\log P(X|Y)\right]\\
    H(X,Y) &=H(Y) + H(X|Y)
\end{align*}

\subsubsection{Chain Rule for $X_1,\ldots,X_n$}
For random variables $X_1,\ldots,X_n$ we have 
\begin{align*}
    H&(X_1,\ldots,X_n)=\\
    &H(X_1) + H(X_2|X_1) + H(X_3|X_1,X_2) + \cdots + H(X_n|X_1,\ldots,X_{n-1})
\end{align*}

\subsection{Subadditivity}
It's just 
\[
H(X),H(Y)\leq H(X,Y)\leq H(X) + H(Y)
\]
Equality in the second inequality if and only if $X,Y$ are independent.

It's very intuitive! The uncertainty in $(X,Y)$ is more than the uncertainty in $X$ or $Y$, separately, but no more thant the sum of uncertainties as we can reveal $X$ and simultaneously reveal $Y$

\subsubsection{Proof}

Follows by:
\begin{align*}
    H(X,Y) &= H(Y) + H(X|Y) \quad \text{and}\quad 0\leq H(X|Y) \leq H(X)\\
    H(X,Y) &= H(X) + H(Y|X) \quad \text{and}\quad 0\leq H(Y|X) \leq H(Y)\\
\end{align*}

The first inequality is quite straight forward. The second one follows from plugging $\leq H(X|Y) \leq H(X)$ into $H(X,Y) = H(Y) + H(X|Y)$.

\subsubsection{Subadditivity for $X_1,\ldots, X_n$}

\[
H(X_1,\ldots,X_n)\leq \sum_i H(X_i)
\]

\subsection{Equalities and Inequalities for Conditional Entropy}

All equalities and inequalities that we have learned (and more generally all equalities and inequalities that involve linear combinations of entropies) generalizes to conditional entropies.

\subsubsection{Examples}

\begin{itemize}
    \item $H(X|Y,Z)\leq H(X|Z)$
    \item $H(X,Y|Z) = H(Y|Z) + H(X|Y,Z)$
    \item $H(X_1,\ldots,X_n|Z)= \cdots$ See photo.
    \item $H(X,Y|Z) \leq H(X|Z) + H(X|Z)$
    \item $H(X_1,\ldots, X_n| Z)\leq \sum_i H(X_i|Z)$
\end{itemize}

\subsubsection{Proof sketch}
To prove 
\begin{itemize}
    \item $H(X|Y,Z)\leq H(X|Z)$
    \item $H(X,Y|Z) = H(Y|Z) + H(X|Y,Z)$
\end{itemize}
It follows by 
\[
H(\cdot | Z) = \sum_c P(Z=c)\cdot H(\cdot | Z=c)
\]
where $\cdot$ stands for any of the above entropies of conditional entropies. So any equality or inequality that holds conditioned on each event $Z=c$ separately also holds when we take the expectation over all possibilities for $Z=c$.

But we need to be careful when there are premises under which the equalities and inequalities are true.

\section{Mutual Information}

\subsection{Setup}

How did Sir Walter Raleigh measure the weight of the smoke of a cigar?

He weighted the cigar before the he smoked it and subtracted the weight of the ashes that remained after he smoked the cigar. The difference was the weight of the smoke.

Scientifically, this is completely wrong.

\subsection{Definition}

$I(X;Y)=$ the information that $X$ gives on $Y$
\[
I(X;Y) = H(Y) - H(Y|X) \geq 0
\]
Remember, this is a definition!

\subsection{Mutual Information is Symmetric}

\begin{align*}
    I(X;Y) &= H(Y) - H(Y|X)\\
    &= H(Y) - \left(H(X,Y)-H(X)\right)\\
    &= H(Y) + H(X) - H(X,Y)\\
    &= I(Y;X) = H(X) - H(X|Y)
\end{align*}
The information that $Y$ gives on $X$ equals to the information that $X$ gives on $Y$. It's a bit counterintuitive, but the above holds and is very useful.

\subsection{Some Simple Facts}

\begin{itemize}
    \item $I(X;Y)=0$ iff $X,Y$ are independent
    \item $I(X;X)=H(X)$
    \item $I(X;f(X))= H(f(X)) \leq H(X)$ for any (deterministic) function $f$
    \item $I(X;Y,Z)\geq I(X;Y),I(X;Z)$
\end{itemize}

\subsection{Conditional Mutual Information}

\begin{align*}
    I(X;Y|Z) &= H(Y|Z) - H(Y|X,Z)\\
    &= H(Y,Z) - H(Z)- \left(H(Y,X,Z) - H(X,Z)\right)\\
    &= H(Y,Z) + H(X,Z) - H(Z) - H(Y,X,Z)\\
    &= I(Y;X|Z)\geq 0
\end{align*}

As before, all equalities and inequalities (that involve linear terms) generalize to the case when conditioning on an additional random variable $Z$, as any such equality/inequality is an equality/inequality about entropies.

\subsection{Chain Rule for Mutual Information}

\begin{align*}
    I(X_1,\ldots,X_n;Y) = I(X_1;Y) &+ I(X_2;Y|X_1) + I(X_3;Y|X_1,X_2)+\\
    &\ldots + I(X_n;Y|X_1,X_2,\ldots, X_{n-1})
\end{align*}
and similarly, we can further condition everything on $Z$.

\section{Concavity of the Entropy Function}

Let $P_1=(P_{1,1},\ldots P_{1,m})$ and $P_2=(P_{2,1},\ldots P_{2,m})$ be two distribution on $m$ elements. Let $0\leq \lambda \leq 1$. Then,
\[
\lambda\cdot H(P_1) + (1-\lambda) \cdot H(P_2)\leq H(\lambda \cdot P_1 + (1-\lambda)\cdot P_2)
\]

\subsection{Proof}

See the photo.

It's essentially just a restatement of the fundamental inequality of conditional entropy $H(X|Y)\leq H(X)$.

\section{Markov Chains}

We say that $X\to Y\to Z$ is a Markov Chain if $P(Z|X,Y)=P(Z|Y)$, that is, for every $a,b,c$ (such that $P(X=a,Y=b)>0$)
\[
P(Z=c|X=a,Y=b)=P(Z=c|Y=b)
\]
Intuitively, $Y$ is generated from $X$ and independent random bits and then $Z$ is generated from $Y$ and independent random bits (without knowing $X$). Each arrow can be viewed as a communication channel or some data processing. $Y$ is a data processing of $X$ and $Z$ is a data processing of $Y$.

By the chain rule for probabilities, we have,
\begin{align*}
    P&(X=a,Y=b,Z=c)=\\
    & P(Y=b)\cdot P(X=a|Y=b)\cdot P(Z=c|X=a,Y=b)
\end{align*}
Thus, $P(Z=c|X=a,Y=b)=P(Z=c|Y=b)$ is equivalent to 
\[
P(X=a,Y=b,Z=c)=P(Y=b)\cdot P(X=a|Y=b)\cdot P(Z=c|Y=b)
\]
for every $a,b,c$. This is an equality, so if we have a Markov chain we get $P(Z=c|X=a,Y=b)=P(Z=c|Y=b)$, and if we have $P(Z=c|X=a,Y=b)=P(Z=c|Y=b)$, we have a Markov chain.

Thus,
\begin{enumerate}
    \item $Z\to Y\to X$ is also a Markov Chain 
    \item Conditioned on $Y$, we have that $X,Z$ are independent
    \item $I(X;Z|Y)=0$
\end{enumerate}

\subsection{Data Processing Inequality}

If $X\to Y\to Z$ is a Markov Chain, then 
\[
I(X;Z) \leq I(X;Y)
\]
or equivalently
\[
H(X|Z) \geq H(X|Y)
\]
The equivalence is because 
\begin{itemize}
    \item $I(X;Z)=H(X)-H(X|Z)$
    \item $I(X;Y)=H(X)-H(X|Y)$
\end{itemize}
The theorem is very intuitive in the data processing interpretation: if we further process the data, we reduce the information that it gives about the source.

\subsubsection{Proof}

If $X\to Y\to Z$ is a Markov Chain, then $Z\to Y\to X$ is also a Markov Chain. We can write 
\begin{align*}
    P(X|Y,Z) &= P(X|Y)\\
    \mathbb{E}\left[\log P(X|Y,Z)\right] &= \mathbb{E}\left[\log P(X|Y)\right]\\
    H(X|Y,Z) &= H(X|Y)
\end{align*}
Hence,
\[
H(X|Z) \geq H(X|Y,Z) = H(X|Y)
\]

\section{Relative Entropy}

Suppose $p=\left(p_1,\ldots, p_m\right)$ and $q=\left(q_1,\ldots, q_m\right)$ are two distributions on $m$ elements. The relative entropy\footnote{We will prove that this is always equal to or greater than 0 via Jensen inequality.} is 
\[
\operatorname{D}\left(p\| q\right)=\sum_i p_i\log \left(\frac{p_i}{q_i}\right)\geq 0
\]
We read it as ``D'' of ``$p$'' relative to ``$q$''.

Relative entropy may be infinite. We can prove the inequality by a convexity argument. When $q$ is uniform, it is the same as Shannon Entropy (up to a sign and an additive constant). We will learn later in the course what $\operatorname{D}\left(p\| q\right)$ really is.

\section{Some exotic entropy functions}

\subsection{Entropy of a Random Variable over the Integers}

Let $X$ be a discrete random variables (possibly infinite number of states), say over the integers 
\[
H(X)=-\sum_{i=-\infty}^{\infty}P(X=i)\log P(X=i)=-\sum_i P_i \log P_i
\]
as in the finite case. The sum may be infinite, but otherwise all equalities and inequalities generalize.

\subsection{Entropy of a Continuous Random Variable}

If $X$ is a continuous random variable, we define 
\[
H(X) = -\int f(x)\log f(x) \, dx
\]
where $f(x)$ is the probability density function (in analogy to the finite case). We assume that $X$ has a ``nice'' probability density function and that the integral is well defined. This is also called Differential Entropy. It may be negative. Equalities generalize, \textit{but inequalities not necessarily}.

\subsubsection{Differential Entropy may be Negative}

Note that $H(X)$ may be infinite, and also negative. For example $f(x)=\frac{1}{a}$ on the interval $[0,a]$. Then
\[
-\int_0^a f(x)\log f(x)\, dx = \log a
\]
which is negative when $a< 1$. This destroys all inequalities (but equalities remain the same).

\subsection{Relative Entropy in the Continuous Case}

If $f,g$ are sufficiently nice probability density functions 
\[
\operatorname{D}(f\| g)=\int f(x)\log\left(\frac{f(x)}{g(x)}\right)\, dx >0
\]

\section{Approximating Binomial Coefficients (informal)}

Let $k\leq n$ be integers. Let $p=\frac{k }{n }$. We have 
\[
\binom{n }{k }=\frac{n!}{k!\cdot \left(n-k\right)!}
\]
Stirling's approximation is $\ln(m!)\approx m\ln m-m$. Thus,
\[
\ln\binom{n }{k }\approx \left(n\ln n -n\right)- \left(k\ln k -k\right) - \left(\left(n-k\right)\ln\left(n-k\right)-\left(n-k\right)\right)
\]
which then can be written as 
\[
\ln\binom{n }{k }\approx n\ln n -k\ln k - \left(n-k\right)\ln \left(n-k\right) = -k \ln \left(\frac{k}{n}\right)\cdots
\]

See the photo.

\subsection{Approximating Multinomial Coefficients}

See the photo.


\end{document}